\section{\texorpdfstring{完备性让 \(L^p\) 成为 Banach 空间}{完备性让 Lp 成为 Banach 空间}}
\label{sec:12-Real-Analysis-Support/07-Lp-Spaces-and-Function-Geometry/02}

你要在一族函数里找一个最好的：使某个泛函最小的那一个。标准动作是造一列 \(f_1,f_2,\dots\)，让目标值一路下降，逼近下确界。序列彼此越靠越近，你写下一句``取极限，记为 \(f^\star\)''，然后开始讨论 \(f^\star\) 的性质。

这一步藏着一张没兑现的支票。序列越靠越近，只说明它是 Cauchy 列；\emph{它并没有说极限存在}，更没有说极限还留在你研究的那一族函数里。有理数列 \(1,1.4,1.41,1.414,\dots\) 也越靠越近，极限却不是有理数。函数空间上会不会重演这一幕？

\subsection{一个逼近序列可以走出自己的家}

会。而且反例很朴素。

在 \(C[0,1]\)（\([0,1]\) 上的连续函数）上装 \(L^1\) 范数 \(\norm{f}_1=\int_0^1|f|\)。取一列连续的斜坡

\[
f_n(x)=
\begin{cases}
0, & 0\le x\le \tfrac12,\\[2pt]
n\left(x-\tfrac12\right), & \tfrac12<x<\tfrac12+\tfrac1n,\\[2pt]
1, & \tfrac12+\tfrac1n\le x\le 1 .
\end{cases}
\]

每个 \(f_n\) 都连续。两项之差只在一小段上非零，且被 \(1\) 界住，所以对 \(m>n\) 有 \(\norm{f_n-f_m}_1\le 1/n\)，这列函数在 \(L^1\) 范数下是 Cauchy 的。

但它逼近的是谁？斜坡越来越陡，极限是阶跃函数 \(f=\ind_{[1/2,1]}\)，而且 \(\norm{f_n-f}_1\le 1/n\to0\)，收敛是货真价实的。阶跃函数不连续，不属于 \(C[0,1]\)。于是 \((C[0,1],\norm{\cdot}_1)\) 有一个 Cauchy 列没有极限——空间上有洞。

值得停一下品品这个洞的性质：洞不在函数上，在\emph{范数}上。同一族 \(C[0,1]\) 配上 sup 范数就是完备的，因为一致收敛保持连续性，上面那列斜坡在 sup 范数下压根不是 Cauchy 列（相邻两项的最大差恒为 \(1\) 的量级）。完备与否是空间和范数配对之后才谈得上的性质，换一把尺子，洞的位置就变了。

\begin{center}
\begin{tikzpicture}[>=stealth]
% 大空间 L^1
\draw[thick] (0,0) ellipse (3.6cm and 2.0cm);
\node at (0,1.65) {\footnotesize \(L^1[0,1]\)：完备，洞被补齐};
% 小空间 C[0,1]
\draw[dashed,thick] (-1.0,-0.1) ellipse (2.0cm and 1.25cm);
\node at (-1.0,-0.55) {\footnotesize \(C[0,1]\) 配 \(L^1\) 范数};
% Cauchy 列
\foreach \x/\y in {-2.55/0.30, -1.95/0.28, -1.35/0.25, -0.75/0.22, -0.25/0.19, 0.15/0.17, 0.45/0.16, 0.65/0.155}
  \fill (\x,\y) circle (0.055);
\node[below] at (-2.55,0.22) {\footnotesize \(f_1\)};
\node[below] at (-1.95,0.20) {\footnotesize \(f_2\)};
\node[below] at (-1.35,0.17) {\footnotesize \(f_3\)};
% 极限点
\draw[fill=white,thick] (1.35,0.15) circle (0.1);
\draw[->,thick] (0.80,0.155) -- (1.20,0.15);
\node[right] at (1.55,0.15) {\footnotesize \(f=\ind_{[1/2,1]}\)};
\node[right] at (1.55,-0.30) {\footnotesize 不连续，落在虚线外};
\end{tikzpicture}
\end{center}

\emph{Cauchy 列在虚线圈内步步逼近，极限却落到圈外；完备性就是禁止这种逃逸。}

\subsection{完备性是所有存在性论证的地基}

有了这张图，完备性的含义可以压成一句：\emph{逼近序列不会逃出空间}。凡是 Cauchy 列，极限就在屋里等着。

为什么这条性质如此关键？因为存在性证明几乎只有一种套路：造一列越来越好的候选，然后取极限。极小化序列、迭代算法的迭代点、级数的部分和、逐步细化的逼近方案，全是这个模式。而取极限这一步，靠的是完备性签字。空间不完备时，你辛辛苦苦造出来的极限可能根本不是你研究的对象——它存在于某个更大的空间里，对你的问题却是废票。

这也解释了实数为什么必须是完备的。\hyperref[sec:12-Real-Analysis-Support/01-Real-Numbers-and-Completeness/01]{``上确界填补有理数的空洞''}讲的是同一件事在最底层的版本，\hyperref[sec:12-Real-Analysis-Support/01-Real-Numbers-and-Completeness/02]{``Cauchy 判据无需预知极限''}讲的是它的实用价值：不必事先知道极限是谁，只看序列自身的收缩就能断言极限存在。\hyperref[sec:12-Real-Analysis-Support/03-Topology-Basics/03]{``完备性让逼近留在空间内''}把这条性质抽象到一般度量空间。本节要说的，是这一整条线终于落到了函数上。

对做数据的人，这条地基的用处相当直接。你在函数空间上写下 \(\argmin_{f\in\mathcal H}\) 某个风险泛函，要断言最小值真的被某个 \(f\) 取到，第一件事就是确认 \(\mathcal H\) 完备（通常再加上闭性与某种紧性或强凸性）。核方法的表示定理、变分问题的可解性、Fourier 展开的收敛，背后都站着同一句话。反过来说，如果有人在一个不完备的函数类上宣称``最优解存在''，那多半是把下确界当成了最小值。

\subsection{Riesz--Fischer 定理把洞补齐}

\begin{theorem}[Riesz--Fischer]
对 \(1\le p\le\infty\)，\(L^p(\mu)\) 在 \(\norm{\cdot}_p\) 下完备，因而是 Banach 空间。
\end{theorem}

证明的骨架是四步。第一步收缩：设 \((f_n)\) 是 Cauchy 列，抽子列使 \(\norm{f_{n_{k+1}}-f_{n_k}}_p<2^{-k}\)，把``任意小''换成一条可求和的几何级数。第二步造包络：令

\[
g=\abs{f_{n_1}}+\sum_{k\ge1}\abs{f_{n_{k+1}}-f_{n_k}},
\]

它的部分和的 \(L^p\) 范数被 \(\norm{f_{n_1}}_p+1\) 一致界住，单调收敛定理让这个界传给 \(g\) 本身，于是 \(g\in L^p\)，特别地 \(g<\infty\) 几乎处处。第三步造极限：\(g\) 几乎处处有限意味着那个级数在几乎每一点上绝对收敛，逐点求和就得到一个函数 \(f\)——极限不是凭空取来的，是用子列一项项造出来的。第四步收网：Fatou 引理给出

\[
\norm{f-f_{n_k}}_p^p
=\int\liminf_{j\to\infty}\abs{f_{n_j}-f_{n_k}}^p\,d\mu
\le\liminf_{j\to\infty}\norm{f_{n_j}-f_{n_k}}_p^p,
\]

右端随 \(k\) 增大趋于零，子列按范数收敛；再用一次 Cauchy 性把收敛从子列传回全列。

对一遍条件的账：第二步的一致界要用 Minkowski 不等式，也就是要用 \(p\ge1\)——上一节那个 \(p=1/2\) 的反例在这里直接掐断整条链。第二、四步分别调用单调收敛与 Fatou，它们都是 Lebesgue 积分的收敛定理，而不是 Riemann 积分能提供的，这是 \(L^p\) 理论必须建在 Lebesgue 框架上的原因之一。第三步的``几乎处处有限''只给出几乎处处的极限，所以造出来的 \(f\) 天然是个等价类，与上一节的约定接得上。最底层，Fatou 与单调收敛又都依赖 \(\R\) 自身的完备性——函数空间的完备性并非凭空而来，是把实数的完备性搬了上来。

\(p=\infty\) 单走一条更短的路：\(\norm{f_n-f_m}_\infty\) 小意味着除去一个零测集后 \(f_n\) 一致 Cauchy，一致极限直接存在且仍被同一个界控制。

\subsection{完备不保证逐点世界也整齐}

完备性管住的是``极限留在空间里''，不管``极限在每个点上表现良好''。\(L^p\) 中的极限仍是等价类，问它在某一点的值依然不合法。需要逐点控制时，要么加强收敛方式（一致收敛），要么换一个空间（连续函数空间、Sobolev 空间、再生核 Hilbert 空间），要么援引嵌入定理把 \(L^p\) 信息升级成逐点信息。

端点上还有一处不对称值得记住：连续函数在 \(L^p\)（\(p<\infty\)）中稠密，所以``先对光滑函数验证、再取极限''是合法策略；但连续函数在 \(L^\infty\) 中\emph{不}稠密，阶跃函数无法被连续函数一致逼近。本节开头那列斜坡正是这句话的注脚——它在 \(L^1\) 意义下补上了洞，在 sup 意义下永远补不上。

\subsection{范数收敛拿不到时，还有弱收敛}

实际问题里常出现这样的局面：序列有界，却怎么也证不出按范数收敛。Hilbert 空间（下一节的 \(L^2\) 就是）提供了退一步的选项：称 \(f_n\) 弱收敛到 \(f\)，若对每个 \(g\) 都有 \(\ip{f_n}{g}\to\ip{f}{g}\)。

范数收敛推出弱收敛（Cauchy--Schwarz 一步到位），反过来不成立。标准反例是 \([0,2\pi]\) 上的 \(f_n(x)=\sin(nx)\)：Riemann--Lebesgue 引理说 \(\int f_n g\to0\) 对一切 \(g\in L^2\) 成立，所以 \(f_n\) 弱收敛到 \(0\)；可是 \(\norm{f_n}_2^2=\pi\) 岿然不动。越来越快的振荡被任何固定的 \(g\) 平均掉了，能量却一点没少。

弱收敛值钱的地方在于它便宜：Hilbert 空间中的有界序列必有弱收敛子列。许多变分与优化论证就靠这一条推进——先取弱极限保证候选存在，再单独验证它满足约束、或者用下半连续性验证它确实达到下界。

到这里，\(L^p\) 已经是一个能安心取极限的场所。可完备性只承诺``极限在屋里''，并不告诉你极限在哪、怎么算。想把``离某个子空间最近的那个函数是谁''算出来，光有范数不够，还需要角度。下一节会看到，只有 \(p=2\) 提供了角度。
